\chapter[Sermon 49. The Enterprises of Antichrist in Weeding Out the Preachers to Be Vain: How Great Shall Be the Rewards of Preachers, and of the Punishment of the Wicked.]{The Enterprises of Antichrist in Weeding Out the Preachers to Be Vain: How Great Shall Be the Rewards of Preachers, and of the Punishment of the Wicked.}
\label{sermon:049}
\markright{Sermon 49. The Enterprises of Antichrist in Weeding Out the ...}

And after three days and a half, the spirit of life from God entered into them. They stood upon their feet, and great fear came upon those who saw them. They heard a great voice from heaven saying to them, “Come up here.” They ascended into heaven in a cloud, and their enemies saw them. In that same hour, there was a great earthquake, and a tenth part of the city fell. In the earthquake, seven thousand names were slain, and the remnant were terrified and gave glory to God of heaven.

\index{Antichrist}Hitherto he has spoken of the wicked joys and gladness of Antichrist and the ungodly men of the last age, concerning the slaughter of the holy prophets of God. They will think how they shall reign forever in those their errors, superstitions, and pleasures; and suppose by their murdering to have silenced the preaching of the gospel, which is most unpleasant to them. But consequently, the Lord shows that their hope is most vain, their attempts to be frustrated, and their joys short; yea, and quickly to be turned into mourning and misery. For first, he declares that the prophecy or preaching shall be repaired by God through new prophets, and that to the greatest grief and terror of the Antichristians, who looked for no such thing. After this, he shows how great rewards are prepared and given to the preachers oppressed in this world and treated with great villainy. Finally, he signifies that the wicked shall not live in continual pleasure, but that God will disturb their joys, bringing misery upon them even in this world. Which, although he begins to punish at the last in this world, will in another world more abundantly augment their everlasting torments. And all these things shall need no great exposition, so that we mark diligently what things have been done a few ages past, and what are done also at this day. And all these things appertain to the consolation and comfort of the Saints.

\index{Erasmus of Rotterdam}\index{Jerome, Saint}\index{Luther, Martin}\index{Oecolampadius, Johannes}\index{Rome}First that the free preaching of God’s word against Antichrist shall be restored, which seemed to him self to have overcome and oppressed all prophecy, he declares by these words: and after three days and an half, the spirit of life from God entered into them. He signifies by that number of days, as I told you before, a very short time, as though he should say: they shall not long enjoy their false and bloody pleasures. For God shall raise up other prophets in the place of those that are dead. And he speaks as though God should raise up the self same prophets, which Antichrist had slain, and that he would object them again to the wicked in their own bodies. How be it they shall be raised again in their bodies at the last day: but now shall other preachers succeed in the place of those that rest, unto whom God shall give that spirit of his, which he had given to the others that are dead. Therefore he calls this the spirit of life, for as much as those which were slain for the same doctrine, seem as it were to have lived again. Verily for likeness of doctrine, John the Baptist, Elijah, and the prophet Jeremiah seemed to have been revived in Christ, as is read in Matthew 14:12 and 16:14. And here is expressly said, that the same spirit did not proceed of the Devil, or of men (as it is said at this day of many), but of God. For he with his spirit (which is one) inspires his ministers, and directs the same by his word, that the latter wholly answer to the former in doctrine, and severe rebuking of sins, and so forth. For the lively effect of that spirit follows, and they stood upon their feet: that is to say, they lived again. Their doctrine seemed overthrown and trodden under foot, but God’s word stands again upon its feet, and runs most swiftly. We say in Dutch of such as be restored, to expound the effect that same also pertains, that the Antichristians saying other preachers succeed in the room of the that were slain: being stricken with fear, know not whether to turn them. By the way therefore is signified, that the course of the word shall be fortunate, and the which these men can not stop by any means, however they rage and murder. All these things shall be better understood by the Histories of later times, and of such things as are done yet at this day. And to the intent, that omitting the eldest things, I may touch those of latter time: the Bishops of Rome had thought they had won the field in the council of Constance, when they had burned John Huss and Jerome of Prague: but within a short time after many godly and well learned men sprang up in Bohemia and in other countries, in whom those slain appeared to have taken again the spirit of life. In Italy Laurence Valla taught to his great praise, and also Jerome Savonarola, and so forth. In Germany taught many godly men, as in France also in Englande, and other nations. Thirty years passed through the grace of God was brought a light into the world by Mirandola, Reuchlin, Erasmus, Luther, Zwinglius, Oecolampadius, Melanchthon, and innumerable others, in whom the spirit of life uttering itself after every man’s talent, set forth the Scriptures, detected the Romish wickedness, and rebuked the vices of all states, but especially of the clergy. The Romans are afraid of this spirit, and fill the ears of the emperor and kings with complaints and accusations, and cry out that we should all with our books be destroyed and burned. How be it the power of God nevertheless makes the prophets to stand on their feet, and their preaching to run a pace: however these rage in their fury, and persecute God’s verity preached throughout the whole world. To God be the praise and glory.

\index{Papacy}\index{Judgment, Last}\index{Daniel (Book of)}In this consolation are mixed also rewards prepared for faithful ministers, whom the Antichristians, after slaying, first excommunicate, so that they may seem bound and delivered to the Devil, to be tormented with everlasting punishment. And until now, all preachers who have spoken against the church of Rome and have suffered therefore at the Pope’s command, have been thought to have perished both in body and soul: their bodies, I say, consumed with fire, and their souls cast down into hell. For they were condemned as heretics and enemies of God and the church, even as plagues of mankind, and so taken out of this life. But contrariwise, the Lord here pronounces and declares everlasting rewards to be prepared for them. For their souls, released from their bodies, are straightway taken up into heaven, and their bodies raised at the last judgment, ascend into heaven also, that there they may rejoice with Christ forevermore. But to the intent that this godly promise of the everlasting and inestimable reward might be of more authority and credit with all men, the Lord propounds it not simply, but most gloriously decked and adorned, for he sets before that a voice was sent to the prophets, and that from heaven; moreover, great and loud. For great is the agreement of Patriarchs, Prophets, and Apostles with the very Son of God in assured doctrine, upon which we believe without doubt that those who suffer for the confession of Christ are saved both in body and soul. And that doctrine was brought from heaven, so that there is no place left for doubtfulness. There are testimonies in the scriptures, both manifest and many, as in Isaiah 26, Daniel 12, Matthew 10 and 16, John 14, and diverse others. What should we say is now brought as an express testimony hereof? For a voice sounds from heaven over the afflicted by the tyranny of Antichrist: “Come up here.” That is to say, as much as, “I see the lewdness and cruelty of the Antichristians to be such that there is no place left for you on earth. They torment and persecute you as plagues, and you are unworthy to live on earth; come therefore here to me, into the heavenly palace, whither I myself came also after the cross and ignominious death.” We read in the Gospel that the judge shall say to the righteous, “Come, you blessed of my father, \&c.”

\index{Resurrection}Furthermore, lest any man should think these words to be vain, the Lord adds through John, “and they ascended into heaven,” not because the resurrection has already occurred, but for the certain knowledge of the thing he speaks of, he speaks of the future as if it were past. Many similar phrases are found throughout the Prophets. Elijah in times past ascended into heaven, both soul and body, as we read in the fourth book of Kings, chapter two, by the same miracle he showed them also what reward the Lord has prepared for faithful preachers of God’s word, and nothing else is repeated here. He adds how they went up in a cloud. For a cloud took up Christ, our head, from the eyes of the disciples, and we also shall be taken up in a cloud to meet the Lord in the air, as the scripture states in the first book of Acts, and in first Thessalonians, chapter four. Although therefore preachers, and those who believe the preachers, are excommunicated by Antichrist, and through open and shameful punishments would seem to be sent to the Devil, Christ receives them, delivered from all evils, unto himself, into the palace of Heaven.

He adds another point as well: their enemies saw them. They saw, I say, with a terrifying fear, for as they witness those whom they condemned as God’s enemies being in glory as true and honorable friends of God, they will conclude that they themselves are destined for fellowship with the Devil. Read a full commentary on this in chapters 3 and 5 of the book of Wisdom. Although, therefore, preachers of the Gospel in this present world are judged and appear to the world as damned, in that day when all people will be assembled—all who ever have been, are now, or will be—it will be manifest to all that these are God’s dearest friends, and that their cause is just. With this, the Lord comforts those who are persecuted, condemned, despised, and scorned for preaching God’s word. By these things, he prepares and establishes the minds of the faithful, so that they are not discouraged by the rebukes, revilings, and oppressions of Antichrist and his followers.

\index{John, Saint}Finally, the Lord also adds certain things concerning the miseries of the Antichristians, with which the righteous Lord begins to punish them and to interrupt their wicked joys, so that at last in another world he may put the same to torment that never shall have end. In that same hour, that is, doubtless in the time wherein they shall afflict the prophets, there shall be a great earthquake, and the tenth part of the city shall fall. And the tenth part we understand to be great, yet so that the greater part shall remain in error. As Saint Peter prophesied should come to pass, in 2 Peter 2, and the Lord himself also in Matthew 7. He seems to recite two evils which hang over them: calamities and revolts. For Saint John himself seems to add an exposition, saying: and there were slain in the earthquake the names of seven thousand men. And the remainder were afraid and gave glory to God of Heaven.

Therefore I suppose the earthquake to signify exceedingly great alterations, commotions, seditions, wars, slaughters, and destructions. And he said the names of men after the Hebrew phrase, for a number of men. And he put 7,000, a certain number, for an uncertain one: as where it is said to Elijah, “I have left me seven thousand men, who have not bowed their knees to Baal.” For if it signifies a great multitude, likewise he signifies here that no small number of Antichristians shall be dispatched out of the way by slaughter and sundry or all kinds of calamities. Again, he signifies that the tenth part of the world, that is to say, the adherents and favorers of the Roman church, will revolt, not a few of them, from the same church, being frightened by the preaching of God’s word, and by plagues inflicted upon the enemies of God’s word, and so they shall forsake the Roman church, that they may give all glory to the God of Heaven.

Having been misled up to this point by Roman trivialities and deceptive opinions, they have not given all glory wholly to the true God, creator of heaven and earth, and the inhabitant and ruler of heaven, while attributing more to creatures, human inventions, and errors than to truth. They have shared the glory that belongs to God alone with saints and the works of their hands. But now, being instructed by the preaching of the gospel, they will depend on God alone and will ascribe all glory to him through Christ.

Now if you consult histories, not ancient (for I would not trouble you with a lengthy recounting), but recently written, within these last hundred years, you will find a remarkable clarity regarding this matter. When the preachers of Bohemia were burned at Constance, a great disturbance of the people immediately arose, the Bohemians waging a deadly war against the Romans. Aeneas Silvius himself wrote of that war, in which many thousands of men were slain, and many places destroyed and laid waste. Moreover, countless people deserted the authority of Rome. In our memory, a grievous persecution was stirred up against the faithful through the instigation of Rome, and thousands of faithful were slain, besides the expectation of all people. Rome was taken in the year of our Lord 1527, and so devastated and plundered that this calamity could be compared with the most severe that ever occurred. Neither do the princes cease to wage war among themselves, and to weaken themselves with mutual destruction, which consistently shed the blood of the faithful. But we are glad and rejoice that a wonderful number at this day revolt from that Roman see, and give to God through Christ all glory. To him be glory and rule forever and ever. Amen.
